% Adapted notation examples for visual regression, not reproductions of figures.
\documentclass{article}
\usepackage[a4paper,margin=20mm]{geometry}
\usepackage{derham}
\usepackage{hyperref}
\begin{document}
\section*{FEEC diagram gallery}
These diagrams are adapted from mathematical relations in the cited sources.
Their labels can be changed without asserting that any particular maps commute.

\subsection*{Conforming finite element projection}
The conforming sequence uses the projection notation of G\"u\c{c}l\"u,
Campos Pinto, and Hadjout, \emph{Journal of Scientific Computing} (2023),
\href{https://doi.org/10.1007/s10915-023-02351-x}{doi:10.1007/s10915-023-02351-x}.
\[
\derhamcommuting[
  dimension=3,
  custom-spaces={V^0,V^1,V^2,V^3},
  custom-discrete-spaces={V_h^{0,c},V_h^{1,c},V_h^{2,c},V_h^{3,c}},
  operators={\operatorname{grad},\operatorname{curl},\operatorname{div}},
  discrete-operators={\operatorname{grad}_h^c,\operatorname{curl}_h^c,\operatorname{div}_h^c},
  maps={\Pi_h^{0,c},\Pi_h^{1,c},\Pi_h^{2,c},\Pi_h^{3,c}},
  column-sep=large
]
\]

\subsection*{Pullback and exterior differentiation}
The reference to physical map follows the pullback convention in
\href{https://www-users.cse.umn.edu/~arnold/feec-cbms/feec.pdf}{Arnold's FEEC notes}.
The diagonal arrow demonstrates a connection between different degrees;
it is a separate, explicitly named map.
\[
\begin{derhamdiagram}[dimension=2,notation=forms,
  custom-spaces={\Lambda^0(\Omega),\Lambda^1(\Omega),\Lambda^2(\Omega)},
  column-sep=large,row-sep=huge]
  \derhamrow[name=physical,kind=continuous]
  \derhamrow[name=reference,kind=reference,
    spaces={\Lambda^0(\widehat\Omega),\Lambda^1(\widehat\Omega),\Lambda^2(\widehat\Omega)}]
  \derhamconnect[from=physical,to=reference,operator=F^*]
  \derhamconnect[from=physical,to=reference,
    from-degree=0,to-degree=1,operator=T]
\end{derhamdiagram}
\]
The additional $T$ arrow is illustrative and carries no implied naturality identity.
\end{document}
